\documentclass[10pt, letterpaper]{extarticle}

% Packages:
\RequirePackage{fontawesome}
\usepackage[
ignoreheadfoot, % set margins without considering header and footer
top=1.25 cm, % seperation between body and page edge from the top
bottom=1.25 cm, % seperation between body and page edge from the bottom
left=1.75 cm, % seperation between body and page edge from the left
right=1.75 cm, % seperation between body and page edge from the right
footskip=1.0 cm, % seperation between body and footer
    % showframe % for debugging 
]{geometry} % for adjusting page geometry
\usepackage{titlesec} % for customizing section titles
\usepackage{tabularx} % for making tables with fixed width columns
\usepackage{array} % tabularx requires this
\usepackage[dvipsnames]{xcolor} % for coloring text
\definecolor{primaryColor}{RGB}{0, 0, 0} % define primary color
\usepackage{enumitem} % for customizing lists
\usepackage{amsmath} % for math
\boldmath
\usepackage[
pdftitle={Hunter Ellis's Resume},
pdfauthor={Hunter Ellis},
pdfcreator={LaTeX with RenderCV},
colorlinks=true,
urlcolor=primaryColor
]{hyperref} % for links, metadata and bookmarks
\usepackage[pscoord]{eso-pic} % for floating text on the page
\usepackage{calc} % for calculating lengths
\usepackage{bookmark} % for bookmarks
\usepackage{lastpage} % for getting the total number of pages
\usepackage{changepage} % for one column entries (adjustwidth environment)
\usepackage{paracol} % for two and three column entries
\usepackage{ifthen} % for conditional statements
\usepackage{needspace} % for avoiding page brake right after the section title
\usepackage{iftex} % check if engine is pdflatex, xetex or luatex
\usepackage{datetime}
\usepackage{lipsum}

% Ensure that generated pdf is machine readable/ATS parsable:
\ifPDFTeX
    \input{glyphtounicode}
    \pdfgentounicode=1
    \usepackage[T1]{fontenc}
    \usepackage[utf8]{inputenc}
    % \usepackage{lmodern}
\fi

% Fonts
\usepackage{libertine}
% \usepackage[sfdefault]{arimo}
% \usepackage{Archivo}

% Some settings:
\raggedright
\AtBeginEnvironment{adjustwidth}{\partopsep0pt} % remove space before adjustwidth environment
\pagestyle{empty} % no header or footer
\setcounter{secnumdepth}{0} % no section numbering
\setlength{\parindent}{0pt} % no indentation
\setlength{\topskip}{0pt} % no top skip
\setlength{\columnsep}{0.15cm} % set column seperation
\pagenumbering{gobble} % no page numbering

\titleformat{\section}{\needspace{4\baselineskip}\bfseries\large}{}{0pt}{}[\vspace{1pt}\titlerule]

\titlespacing{\section}{
    % left space:
    -1pt
}{
    % top space:
    0.15 cm
}{
    % bottom space:
    0.15 cm
} % section title spacing

\renewcommand\labelitemi{$\vcenter{\hbox{\tiny$\bullet$}}$} % custom bullet points
\newenvironment{highlights}{
    \begin{itemize}[
        topsep=0 pt,
        parsep=0 pt, 
        partopsep=0pt,
        itemsep=0pt,
        leftmargin=0.25 cm + 10pt
        ]
    }{
\end{itemize}
} % new environment for highlights


\newenvironment{highlightsforbulletentries}{
    \begin{itemize}[
        topsep=0.10 cm,
        parsep=0.10 cm,
        partopsep=0pt,
        itemsep=0pt,
        leftmargin=10pt
        ]
    }{
\end{itemize}
} % new environment for highlights for bullet entries

\newenvironment{onecolentry}{
    \begin{adjustwidth}{
            0 cm + 0.00001 cm
        }{
            0 cm + 0.00001 cm
        }
    }{
    \end{adjustwidth}
} % new environment for one column entries

\newenvironment{twocolentry}[2][]{
    \onecolentry
    \def\secondColumn{#2}
    \setcolumnwidth{\fill, 4.5 cm}
    \begin{paracol}{2}
    }{
        \switchcolumn \raggedleft \secondColumn
    \end{paracol}
    \endonecolentry
} % new environment for two column entries

\newenvironment{threecolentry}[3][]{
    \onecolentry
    \def\thirdColumn{#3}
    \setcolumnwidth{1.5\fill, 3\fill, 1.5\fill}
    \begin{paracol}{3}
        {\raggedright #2} \switchcolumn
        \begin{center}
        }{
        \end{center}
        \switchcolumn \raggedleft \thirdColumn
    \end{paracol}
    \endonecolentry
}

\newenvironment{header}{
    \setlength{\topsep}{0pt}\par\kern\topsep\linespread{1.5}
}{
    \par\kern\topsep
} % new environment for the header

% save the original href command in a new command:
\let\hrefWithoutArrow\href

\begin{document}
\newcommand{\AND}{\unskip
    \cleaders\copy\ANDbox\hskip\wd\ANDbox
    \ignorespaces
}
\newsavebox\ANDbox
\sbox\ANDbox{$|$}
\LARGE{\textbf{{Hunter Ellis}}}\\
\kern 2.5 pt%

\normalsize
 \hrefWithoutArrow{tel:+1-703-953-6963}{(703) 953-6963} \ |
 \ Atlanta, GA \ |
 \ \hrefWithoutArrow{mailto:elliswhunter@gmail.com}{elliswhunter@gmail.com} \ |
 \ \hrefWithoutArrow{https://github.com/hjkellis}{{github.com/hjkellis}} \ |
 \ \hrefWithoutArrow{https://www.linkedin.com/in/ellishw/}{linkedin.com/in/ellishw} \ |
 \ \hrefWithoutArrow{https://ellishw.tech}{{hjkellis.github.io}}
 \normalsize

\kern -15.0 pt%
\section{}
Engineer focused on building intelligent, capable platforms. Experience in controls, motion planning, estimation, and tracking.\\
% focusing on path planning, control stystems, and sensor fusion

\section{Experience}
\begin{twocolentry}{\textbf{Aug 2025 -- Present} \\\textit{Atlanta, GA}}
    \textbf{Avionics Research Engineer}\\
    \textit{Georgia Tech Research Institute $\cdot$ Applied Systems Lab}\\
\end{twocolentry}
\kern 2.0 pt%
\begin{onecolentry}
    \begin{highlights}
        \item Developing C++ software for collaborative autonomous UAV swarms in a component-based software engineering environment. % Git GNU/Linux and CMake C++14
        \item Integrating, developing, and benchmarking sampling-based motion planning algorithms over constraint manifolds leveraging The Open Motion Planning Library (OMPL). % may be too wordy
        \item Developing a containerized motion planning service with hybrid planning using graph, sampling, and optimization methods.
        \item Leading an independent research project developing and delivering reactive obstacle avoidance.
        \item Implemented multi-hypothesis tracking, including prediction and data association algorithms for multi-object tracking (MOT) in decentralized swarms using distributed camera data.
        % \item Embedded software development
        \item Software systems integration
        \item Refactored legacy code to use DDS middleware.
    \end{highlights}
\end{onecolentry}
\kern 5.0 pt%

\begin{twocolentry}{\textbf{Aug 2023 -- May 2025} \\\textit{Blacksburg, VA}}
    \textbf{Robotics Research / M.S. Computer Engineering}\\
    \textit{Virginia Tech $\cdot$ Control, Optimization, and Online Learning for Autonomy Lab}\\
\end{twocolentry}
\kern 2.0 pt%
\begin{onecolentry}
    \begin{highlights}
    \item Built a 6-DOF robotic arm and an accompanying ROS2–Gazebo simulation environment for training, evaluating, and deploying custom control algorithms.
    \item Integrated object detection (YOLOv8), natural language processing, symbolic reasoning, and a DDPG-based reinforcement learning policy for robotic manipulation tasks.
    \item Aided professors in teaching fundamental concepts in linear systems theory and digital signal processing, including Laplace Transforms, Z-Transforms, system stability, and FIR \& IIR filter design.
            % \item Assisted with hands-on projects to illustrate and integrate analog and digital filter design and application on breadboards and TI MSP432 development boards.
    \end{highlights}
\end{onecolentry}
\kern 5.0 pt%

\begin{twocolentry}{\textbf{May 2024 -- Aug 2024} \\\textit{Huntsville, AL}}
    \textbf{Thrust Vector Control Intern}\\
    \textit{Jacobs Space Exploration Group $\cdot$ Thrust Vector Control Test Lab}\\
\end{twocolentry}
\kern 2.0 pt%
\begin{onecolentry}
    \begin{highlights}
    \item Developed thrust vector control testing hardware and software for NASA's Active Inertial Load Simulator at the Marshall Space Flight Center.
    \item Created and ran tests to develop a mathematical model of an electro-mechanical actuator -- used Python, MATLAB, and LabVIEW.
            % \item Derived control algorithms for a load-simulating actuator, in Simulink, to simulate external loads placed on the thrust vector control actuators during flight. 
            % \item Integrated IIR filters to reduce high frequency noise from a load cell and linear variable differential transformer (LVDT).
    \end{highlights}
\end{onecolentry}
\kern 5.0 pt%

\begin{twocolentry}{\textbf{Jun 2023 -- Aug 2023}\\ \textit{Grenoble, France}}
    \textbf{Control Research Intern}\\
    \textit{G2ELab $\cdot$ Power Electronics Lab}\\
\end{twocolentry}
\kern 2.0 pt%
\begin{onecolentry}
    \begin{highlights}
    \item Implemented different inverter control methods for 4-leg (microgrid) topologies and tested them in Simulink and HIL.
            % \item Simulated neutral point capacitive and balancing topologies using 4-leg inverters in Simulink. Tested PI control, PR control, Clarke and Park Transforms with HIL simulations.
    \end{highlights}
\end{onecolentry}
\kern 5.0 pt%

\begin{twocolentry}{\textbf{Jun 2022 -- Aug 2022}\\\textit{Bethesda, MD}}
    \textbf{Naval Research Enterprise Internship Program}\\
    \textit{Naval Surface Warfare Center $\cdot$ Carderock Division}\\
\end{twocolentry}
\kern 2.0 pt%
\begin{onecolentry}
    \begin{highlights}
    \item Developed concept hospital sea-train designs, estimated fuel consumption, and electrical power loads.
    \end{highlights}
\end{onecolentry}

\section{Skills}
\textbf{Languages: } C++, Python, MATLAB/Simulink, C, Lua, \LaTeX \\
\kern 3.0 pt%
\textbf{Tools \& Libraries: } ROS2, OMPL, DDS, Qt, Simulink, Docker, Git, FreeRTOS, SolidWorks, Autodesk Inventor (Certified)\\
% \kern 3.0 pt%
% \textbf{Domain Knowledge: } multi-object tracking, sampling-based motion planning, Kalman filters, factor graphs \\

% \section{Projects}
%
% \begin{twocolentry}{}
% \hrefWithoutArrow{https://github.com/hunterwellis/Manipulator-Environment}{\textbf{SLAM FPV Drone Simulator} \faExternalLink}
%     \begin{highlights}
%     \item FPV Drone simulator for training FPV pilots, written in C++ with the Godot game engine.
%     \item Uses SLAM algorithms to create custom environments for pilots to train in.
%     \end{highlights}
% \end{twocolentry}
% \kern 2.0 pt%
%
% \begin{twocolentry}{}
% \hrefWithoutArrow{https://github.com/hunterwellis/Manipulator-Environment}{\textbf{6-DOF Robot Arm} \faExternalLink}
%     \begin{highlights}
%     \item 3D printed robot arm, built using stepper motors and pulleys.
%     \item Implemented ROS2 Jazzy control and Gazebo Harmonic simulation.
%     \end{highlights}
% \end{twocolentry}
% \kern 5.0 pt%
%
% \kern 5.0 pt%
% \hrefWithoutArrow{https://ellishw.tech/projects/rocket-lqi/}{\textbf{LQI Rocket Landing Simulation} \faExternalLink}
% \begin{onecolentry}
%     \begin{highlights}
%     \item Landing a \textit{very simplified} simulated rocket in MATLAB using optimal control.
%     \item Designed a LQI controller for full-state feedback and setpoint tracking of a landing trajectory.
%     \end{highlights}
% \end{onecolentry}
% \kern 5.0 pt%

\section{Education}

\begin{twocolentry}{\textbf{Jun 2026}\\\textit{Atlanta, GA}}
    \textbf{Professional Education} \\
    \textit{Georgia Tech $\cdot$ College of Computing}
\end{twocolentry}
\kern 5.0 pt%

\begin{twocolentry}{\textbf{May 2025}\\\textit{Blacksburg, VA}}
    % \textbf{Master of Science in Computer Engineering} (Accelerated Master's Program)\\
    \textbf{M.S. in Computer Engineering} \\
    \textit{Virginia Tech $\cdot$ Bradley Department of Electrical \& Computer Engineering}
    % \quad\quad\textit{Advisors: }\hrefWithoutArrow{https://coolautonomylab.github.io/members/thinh.html}{{\textit{Dr. Thinh Doan (UT Austin)}}} and \hrefWithoutArrow{https://filebox.ece.vt.edu/~mhsiao/}{{\textit{Dr. Michael Hsiao (Virginia Tech)}}}\\
\end{twocolentry}
\begin{onecolentry}
    \begin{highlights}
        \item Focus: Software and Machine Intelligence
        \item Research Group: Control, Optimization, and Online Learning for Autonomy Lab at UT Austin
    \end{highlights}
\end{onecolentry}
\kern 5.0 pt%

\begin{twocolentry}{\textbf{May 2024}\\\textit{Blacksburg, VA}}
    % \textbf{Bachelor of Science in Electrical \& Computer Engineering} (double major)\\
    \textbf{B.S. in Electrical Engineering and Computer Engineering (Double Major)} \\
    \textit{Virginia Tech $\cdot$ Bradley Department of Electrical \& Computer Engineering}
\end{twocolentry}
\begin{onecolentry}
    \begin{highlights}
        \item Focus: Controls, Robotics, and Autonomy
        \item Design Teams: Solar Car and Human Powered Submarine
    \end{highlights}
\end{onecolentry}
\kern 2.0 pt%

% \hrefWithoutArrow{https://ellishw.tech/projects/closed-loop-stepper/}{\textbf{Closed Loop Stepper Motor} \faExternalLink}
%   \begin{onecolentry}
%         \begin{highlights}
%         \item Backdrivable stepper motor driver using closed-loop control and a magnetic encoder for feedback.
%         \item 4-layer PCB mounts to the back of the motor with CAN and power connections.
%         \end{highlights}
%     \end{onecolentry}
% \end{twocolentry}
% \kern 5.0 pt%

% \hrefWithoutArrow{https://solarcaratvt.org}{\textbf{Design Teams | Solar Car \& Human Powered Submarine} \faExternalLink}
% \begin{onecolentry}
%     \begin{highlights}
%     \item Overall E/E architecture of the Solar Car.
%     \item Single board computer and LCD to display relevant data to the submarine pilot.
%     \end{highlights}
% \end{onecolentry}

% \section{Papers and Research}
% \textbf{A Neuro-Symbolic Reinforcement Learning Architecture: Integrating Perception, Reasoning, and Control} -- Virginia Tech\\
% \textit{Master's Thesis, Jun 2025}

\end{document}
